\scanpage{7}
Pour écrire les opérations de $K^\bullet(X)$ sur $K_\bullet(X)$
\inserted{et l'élément \deleted{\unsure{e}} $1$ de $K_\bullet(X)$}, on
\unsure{pourra} supposer (pour simplifier l'écriture)
$X$ \uline{connexe}.

\deleted{$\unsure{K}$} $1^\circ)$ $\dim X=0$
\diagram{\begin{tikzcd}
K^\bullet(X) \arrow[r,"\operatorname{rg}","\simeq"'] & Z\\
K_\bullet(X) \arrow[r,"\operatorname{lg}"] & Z
\end{tikzcd}}
\margin{isom de $\lambda$-anneaux}

Les opérations de $K^\bullet(X)$ sur $K_\bullet(X)$ : \unsure{pas} le choix !

$1\in K_\bullet(X)$ donné par
\[
\operatorname{lg}(1)=\begin{tikzpicture}[baseline=(mu.base)]
\node[draw,ellipse,inner sep=3pt] (all) {$\phantom{\mu_i\in Z}$};
\node at (all.center) {$\phantom{\mu_i}\in Z$};
\node[draw,ellipse,inner sep=2pt] (mu) at ([xshift=-.38cm]all.center) {$\mu_i$};
\end{tikzpicture}
\]
seul invariant de la situation

ici $c^\bullet=\operatorname{ch}^\bullet=\operatorname{id}_{K^\bullet}$
\deleted{$\unsure{\mu=\operatorname{lg}}\unsure{1}$}
$c_\bullet=\operatorname{ch}_\bullet=\operatorname{id}_{K_\bullet}$

$2^\circ)$ $\dim X=1$. On a les invariants suivants

\unsure{c)} L'ens fini $\kappa=\kappa(X)$ des comp. irréd. de $X$

\unsure{e)} La famille $(\mu_i)_{i\in\kappa}$ (des $\operatorname{lg}_i(\underline O_X)$)
\inserted{$\delta\in K_0\unsure{(X)}$\qquad $(=1-\sum\mu_i1_i)$}
\margin{\unsure{un vecteur}, \unsure{dépendant du} splitting choisi
la classe résiduelle de l'élément $1$ de $K_\bullet(X)$}

$a)$ Le groupe $K^{(1)}=\operatorname{Pic}(X)$

$b)$ Le groupe $\unsure{K_{(0)}^\kappa|K_{(0)}}=K_{(0)}(X)$
\deleted{$\unsure{i\in\kappa}\ [\unsure{K_{(0)}=\prod_{i\in\kappa}K_{(0)}^i}$}

$d)$ L'\deleted{\unsure{es}}hom
\diagram{\begin{tikzcd}
u_i:K^{(1)} \arrow[r] & \unsure{K_{(0)}^\kappa|K_{(0)}}
\end{tikzcd}}
\deleted{$\unsure{i\in\kappa}$} $i\in\kappa$

(savoir \deleted{$u:K^{(1)}\longrightarrow\unsure{K_{(0)}}$}
$L\longmapsto\operatorname{cl}((\underline L-1)\underline O_{X_i})$]

Avec cela, on écrit
\diagram{\begin{tikzcd}[column sep=large]
K^\bullet(X)\simeq Z\times K^{(1)} \arrow[r,"c=\operatorname{ch}","=\operatorname{id}"'] & B^\bullet(X)\simeq Z\times K^{(1)}\\
K_\bullet(X)\simeq Z^\kappa\times K_{(0)} \arrow[r,equal,"\deleted{\unsure{\longrightarrow}}"] & B_\bullet(X)\simeq Z^\kappa\times K_{(0)}
\end{tikzcd}}
\margin{isom ``naturel'' [compatible avec changement de base]}
\margin{$K^{(1)}$ : carré nul}
\margin{isom \uline{non naturel}\\
les $\unsure{\underline O_{X_i^{\mathrm{red}}}}$ ne sont pas inv. par changement de base.}

Opérations de $K^\bullet(X)$ sur $K_\bullet(X)$ : il suffit pour
$\forall\ \xi\in K^{(1)}$ de définir dans $K_\bullet(X)$ une opération
\deleted{de carré nul, \unsure{d'élevant la} filtration \unsure{de} $1$}
de degré filtrant $\geq1$ [donc de carré nul]
i.e. un hom $Z^\kappa\longrightarrow K_{(0)}$
\deleted{$\unsure{\prod_{i\in\kappa}K_{(0)}^i}$}, on prend
\[
(l_i)_{i\in\kappa}\longmapsto(l_i u_i(\xi))_{i\in\kappa}
\]

\margin{6}
